% SPDX-License-Identifier: CC-BY-4.0
% Copyright (c) 2026 Martin Schröder <info@swedishembedded.com>

\section{Results so far}\label{sec:results}

\paragraph{First sound of the speaking agent (C6, partial).}
\Cref{tab:latency} shows \LatencyTurns{} spoken questions answered by the cascade of \cref{sec:system} (Qwen3-8B, synthetic voice, seed 1) on
\LatencyGpu{} with the CUDA backend, run \code{\LatencyRunId}. The first sound arrives \FirstSoundPfifty{}\,s after the end of the question at the median and
\FirstSoundPninetyfive{}\,s at the 95th percentile, against a plan of 1.5\,s and 2.0\,s. This is one run in one process on synthetic speech; the gate
requires at least three repeats on a card shown idle, with intervals, and a real-speech subset, so the claim stays partial.

\begin{table}[t]
\centering\small
\input{generated/latency_table}
\caption{Stages of a spoken turn, run \code{\LatencyRunId}. Whole-turn time is the card's cost, not the listener's wait: the audio is longer than the turn.}
\label{tab:latency}
\end{table}

\paragraph{Teaching: facts, rules and procedures (T2).}
\Cref{tab:teaching} is the test split of \cref{sec:evaluation}, run once under \code{prereg/teaching-test.md} (run \code{teaching-test-20261010};
\TeachScenarios{} scenarios, \TeachAsks{} asks; Qwen3-8B, non-thinking, CUDA). Untaught, the model answered none of the questions about the invented facts
(\TeachFactUntaught{}); with the ledger of a session that had been taught, in a new session, it answered \TeachFactTaught{} as taught. A fact that was taught and then corrected was
answered as corrected and never as first taught, a fact taught and then forgotten was not told, and unrelated questions kept their answers (the last two rows are leak and disturbance checks: both arms should pass, and they do). Sixteen facts taught in one
session were each recalled from the ledger.

Three results do not meet the pre-registered gate. The procedure kind has only \(6\) asks, so its interval (lower bound \TeachProcedureLow{}) cannot reach 0.80 even when every ask passes: the gate is not met, and the claim for procedures needs more scenarios, which
would be a new pre-registered set. The rule kind has eight asks and the same limit. And two rules taught together were not both obeyed: every reply ended with ``over'' but none used ``captain'', which the same rule taught alone had
produced in every reply. That is a measured failure of the multi-rule case in one scenario. The agent restated every teaching utterance for confirmation (\TeachUnderstood{} of \TeachTeachings{}), but these utterances were written in the forms the plain reader knows, so that number tests the
protocol, not the reader.

\begin{table}[t]
\centering\small
\input{generated/teaching_table}
\caption{Test split, untaught against taught. Asks that pass their text check, with the 95\% Wilson interval of the taught rate. For the last two rows a pass is the absence of a leak (a forgotten fact told) and of a disturbance (an unrelated answer changed).}
\label{tab:teaching}
\end{table}

\paragraph{Taking a teaching back inside the conversation (T2, addendum).}
A rule that the user has taken back kept being followed in the same conversation, because the model followed its own earlier replies (two dev scenarios failed on all six asks), and a note telling the model the rule was cancelled
made it worse, since it named the rule. The agent now starts the conversation over when a teaching is taken back. Under \code{prereg/teaching-retract.md}, five scenarios (one rule, four facts), asked in the same session after ``Forget that'',
passed \RetractPassed{} asks (run \code{teaching-retract-20261010}). The control arm, which was taught nothing, passes the same absence checks, so this shows that the taught content did not come back, not that anything was learned. The price is that
the conversation so far is gone from the model's view.

\paragraph{Teaching by voice (T4, partial).}
Three scripted conversations were run end to end: the user's turns rendered in one synthetic voice and recognised, the agent hearing only the recognised text, its replies spoken. Under \code{prereg/teaching-spoken.md}, \SpokenPassed{} met their pre-registered conditions:
a corrected fact was restated, confirmed and recalled; a rule was restated, confirmed, obeyed, taken back by ``Forget that'', and not obeyed after; a procedure was restated and told. The recogniser turned ``atlas seven'' into ``\SpokenHeardServer{}'' and ``quartz'' into
``court's''; the agent restated what it heard, which is where the user would catch it, and kept what it heard. These are three conversations in one voice, a demonstration that the path works and not a rate; the first runs of the same conversations found two defects (\cref{sec:system}).

\paragraph{A defect the dev split found.}
The first dev run gave the taught arm exactly the answers of the untaught one: everything was kept, nothing reached the model. The live addendum had been applied outside the
wrapper that sets the system turn, which then replaced it. It was fixed with a test that fails on the wrong order and passes on the right one; the second dev run, after the fix, passed every ask. Neither dev run is a result.

\paragraph{Reading natural phrasings (T1, T3).}
\Cref{tab:reading} scores the readers on the test split of \(121\) authored utterances (\(80\) in test; run \code{teaching-read-20261010}). The plain reader never fired on a non-teaching and found \ReadPlainFound{} teachings;
the model reader found \ReadModelFound{} and fired on \ReadModelFired{} non-teachings. The plain reader's misses are the point of the model reader: ``you got that wrong, the build server is atlas'', ``keep your answers short from now on'', ``to rotate the keys: freeze deploys, mint the new pair, then retire the old pair'' are teachings in
forms its openers do not list. One development false fire of the plain reader (``Next time I will bring coffee'') was fixed before the test split was scored. The utterances are written by the implementer, so these numbers measure coverage of phrasings we thought of, not of the ones users produce; that is the first thing a real corpus must change.

\begin{table}[t]
\centering\small
\input{generated/reading_table}
\caption{Readers of what a user says, on the test split of the authored utterances. ``Misread'' is a reading of a teaching with the wrong kind or without its key content.}
\label{tab:reading}
\end{table}

\paragraph{The recogniser as a judge of pronunciation (a screen, not a result).}
To find words a pronunciation experiment could use, \PilotSentences{} sentences (53 names, each in two carrier sentences) were spoken and recognised by one
recogniser (run \code{pilot-names-20261009}). \PilotDiffering{} came back different from the text after normalisation, but \PilotEmpty{} of those were empty transcripts
and others were cut off or repeated their last words; these are artefacts of one-shot recognition, not mispronunciation. Genuine candidates also appear
(for example Faneuil heard as ``phoneil'', Rittenhouse as ``written house''). The consequence for the design is in \cref{sec:evaluation}: a pronunciation claim needs
two agreeing recognisers and a listener.

\paragraph{A tail of silence before recognition (a screen).}
Of \TakesN{} takes in one pinned voice (run \code{takes-pilot-20261009}), the streaming transducer returned a transcript differing from the text for
\FirstRecogniserRaw{} when the take ended with the speech and for \FirstRecogniserTailed{} when 0.6\,s of silence followed it; the fixed-window
recogniser differed for \SecondRecogniserRaw{} and \SecondRecogniserTailed{}. Both recognisers, given the tail, heard within the word-error bound
\TakesAgreed{} of \TakesN{} takes; they disagreed on \TakesContested{}. This motivates the two-recogniser, tail-padded judge of
\cref{sec:evaluation}; twenty takes are a screen, not a result.

\paragraph{The training machinery learns a planted reward (a gate, not C4).}
Before any claim about naturalness can be tested, the loop that trains the synthesizer from scores of its own takes has to be shown to
work. With a reward that has a known answer (fewer codec frames per word, i.e.\ speaking faster), \GateUpdates{} updates over groups of \GateGroup{} takes moved the
held-out sentences from \GateFramesBefore{} to \GateFramesAfter{} frames per word while the word error rate of the held-out takes, heard by the
streaming recogniser, went from \GateWerBefore{} to \GateWerAfter{}; the KL estimate to the original policy reached \GateKlLast{} at the last update
(run \code{rl-speed-gate-20261010}; eight held-out sentences, three takes each, one run and one seed). This shows the policy is trainable and
generalizes across sentences; it also shows how a reward is exploited, which is why the reward of \cref{sec:formalism} carries a word-error term and the
evaluation a human preference test.

\Cref{tab:methods} puts the same gate under the two baselines the claim has to beat, with the same sentences, takes per update and
seed: imitating the best take of each group moved the held-out sentences least, and the group-advantage update most, with the pairwise update between them
(runs \code{rl-speed-dpo-20261010}, \code{rl-speed-best-of-n-20261010}). On a planted reward this orders the methods by how much of the group they use; it says nothing yet about naturalness,
and one seed on eight held-out sentences cannot separate methods that are close.

\begin{table}[t]
\centering\small
\input{generated/methods_table}
\caption{One planted reward (fewer codec frames per word), three training methods, held-out sentences before and after twelve updates.}
\label{tab:methods}
\end{table}

\paragraph{Reading what a user says (C1, a baseline).}
The rule-based readers of \cref{sec:formalism} were run on \ExtractTestSize{} authored utterances of the test split (the utterances were written, with the
gold label of what the user meant, before the readers were run, and the readers were not changed after; run \code{extraction-rules-20261010}). \Cref{tab:extraction} shows
that the readers fired on no utterance they should not have (\ExtractFalsePositives{} false positives, near-misses such as ``Do you pronounce Jefferson correctly?'' included) and
missed a share of the ways people phrase a correction that differs by reader (the recall column). That is the intended trade: a false lesson is worse than a missed one, and a
missed one costs a repetition. These utterances are written by the implementer, so the table measures conformance to the specification, not accuracy on users;
it is the baseline a model-assisted reader has to beat on utterances written after it is frozen.

\begin{table}[t]
\centering\small
\input{generated/extraction_table}
\caption{Rule-based readers on the test split, each asked only about the turns it would be listening to. Precision and recall of firing on the utterances of its class.}
\label{tab:extraction}
\end{table}

\paragraph{Speech understanding without a transcript (C8, failed).}
A projector trained to put the recogniser's encoder features into the language model's embedding space produced fluent persona answers that were often
about the wrong topic: on 12 held-out questions the word F$_1$ against the text path's answer was 0.299 from speech against 0.489 for the cascade, and the answer named the
topic in 4 of 12 against 9 of 12 (repository roadmap, M4). One voice and 25 templated topics; the gate was not met and no redesign has been trained. Speech-in is
excluded from this paper's claims unless a redesign passes a gate fixed before it is trained.

\paragraph{Everything else.}
C1--C5 and C7 have no result: \unmeasured.
